% SEGMENT OLP-0622-S01
% Part: history
% Chapter: biographies
% Section: georg-cantor

\documentclass[../../../include/open-logic-section]{subfiles}

\begin{document}

\olfileid{his}{bio}{can}

\olsection{ஜியோர்க் காண்டோர்}

\olphoto{cantor-georg}{ஜியோர்க் காண்டோர்}

ஜியோர்க் காண்டோர் (\textsc{gay}-org \textsc{kahn}-tor) பற்றிய தொடக்ககால
வாழ்க்கை வரலாறு ஒன்று, ரஷ்யாவின் செயின்ட் பீட்டர்ஸ்பர்க்கை நோக்கிச் சென்ற
கப்பலில் அவர் பிறந்து கண்டெடுக்கப்பட்டதாகவும், அவருடைய பெற்றோர் யாரென்று
தெரியாததாகவும் கூறியது. ஆனால் அது உண்மையல்ல; எனினும் அவர் 1845-இல்
செயின்ட் பீட்டர்ஸ்பர்க்கில் பிறந்தார்.

% SEGMENT OLP-0622-S02
காண்டோர் 1867-இல் பெர்லின் பல்கலைக்கழகத்தில் கணித முனைவர் பட்டம் பெற்றார்.
கணக்கோட்பாட்டில் செய்த பணிக்காக அவர் அறியப்படுகிறார்; கணக்கோட்பாட்டை ஒரு
தனித்த ஆராய்ச்சித் துறையாக நிறுவியவர் என்ற பெருமையும் அவருக்குரியது.
வெவ்வேறு அளவுகளையுடைய முடிவிலாக் கணங்கள் இருப்பதை முதலில் நிறுவியவர் அவரே.
அவருடைய கோட்பாடுகள், குறிப்பாக முடிவிலிகள் பற்றிய கோட்பாடு, அக்காலக்
கணிதவியலாளர்களிடையே பெரும் விவாதத்தை ஏற்படுத்தின; அவருடைய பணி
சர்ச்சைக்குரியதாக இருந்தது.

% SEGMENT OLP-0622-S03
காண்டோரின் சமய நம்பிக்கைகளும் கணிதப் பணியும் பிரிக்கவியலாத அளவு
பிணைந்திருந்தன; எல்லையற்ற எண்களின் கோட்பாட்டைக் கடவுள் தமக்கு நேரடியாகத்
தெரிவித்ததாகக்கூட அவர் கூறினார். பிற்காலத்தில் காண்டோர் மனநோயால்
பாதிக்கப்பட்டார். 1894 முதல், பின்னாளில் இன்னும் அடிக்கடி, அவர் மருத்துவமனையில்
சேர்க்கப்பட்டார். கணிதவியலாளர் லியோபோல்ட் குரோனெக்கருடனான உறவு முறிவும்
உட்பட, அவருடைய பணிக்கு எழுந்த கடுமையான விமர்சனம் மனச்சோர்வுக்கும் கணிதத்தில்
ஆர்வம் குறைவதற்கும் வழிவகுத்தது. மனச்சோர்வு ஏற்பட்ட காலங்களில் அவர்
மெய்யியல், இலக்கியம் ஆகியவற்றில் ஈடுபட்டார்; ஷேக்ஸ்பியரின் நாடகங்களை
எழுதியவர் பிரான்சிஸ் பேக்கன் என்ற கோட்பாட்டைக்கூட வெளியிட்டார்.

% SEGMENT OLP-0622-S04
காண்டோர் 1918 ஜனவரி 6 அன்று ஹாலேயில் உள்ள ஒரு சிகிச்சை இல்லத்தில் இறந்தார்.

% SEGMENT OLP-0622-S05
\begin{reading}
காண்டோரின் விரிவான வாழ்க்கை வரலாறுகளுக்கு \citet{Dauben1990},
\citet{Grattan-Guinness1971} ஆகியவற்றைப் பார்க்கவும். அவருடைய வழக்கத்திற்கு
மாறான கருத்துகள், பிபிசி ரேடியோ 4-இன் \emph{A Brief History of
  Mathematics} நிகழ்ச்சியிலும் விவரிக்கப்படுகின்றன \citep{Sautoy2014}.
காண்டோரின் கோட்பாடுகளை ராப் இசை வடிவில் கேட்க விரும்பினால்,
\citet{Rose2012}-ஐப் பார்க்கவும்.
\end{reading}

\end{document}
